\chapter{Acetaminophen Level}
\section*{What Is This Test?}
The acetaminophen level measures the serum concentration of acetaminophen (paracetamol, N-acetyl-para-aminophenol, APAP), a widely used analgesic and antipyretic. In overdose, acetaminophen causes severe, potentially fatal centrilobular hepatocellular necrosis. The mechanism of toxicity: at therapeutic doses, acetaminophen is primarily metabolized by glucuronidation and sulfation to non-toxic metabolites. A small fraction (5--10\%) is oxidized by CYP2E1 (and CYP1A2, CYP3A4) to the toxic reactive intermediate N-acetyl-p-benzoquinone imine (NAPQI). NAPQI is normally detoxified by conjugation with glutathione (GSH) and excreted as mercapturic acid and cysteine conjugates in the urine. In overdose (ingestion >7.5 g in adults, >150 mg/kg in children): (1) the glucuronidation and sulfation pathways are saturated, shunting more acetaminophen toward CYP2E1 oxidation and increasing NAPQI production; (2) hepatic glutathione stores are depleted (<30\% of normal), leaving NAPQI unopposed to bind covalently to hepatocyte macromolecules (particularly mitochondrial proteins), causing mitochondrial dysfunction, oxidative stress, hepatocellular necrosis, and acute liver failure. The clinical course of acetaminophen poisoning has four stages: Stage I (0--24 hours): nausea, vomiting, anorexia, malaise, diaphoresis; LFTs normal. Stage II (24--72 hours): right upper quadrant pain, tender hepatomegaly, rising AST/ALT (the first laboratory sign; ALT >1,000 IU/L is characteristic), prolonged PT/INR. Stage III (72--96 hours): peak hepatotoxicity (massive AST/ALT elevations up to 10,000--50,000+ IU/L), fulminant hepatic failure with jaundice, coagulopathy, hepatic encephalopathy, hypoglycemia, lactic acidosis, acute kidney injury (ATN from NAPQI in the kidney or hepatorenal syndrome), and death. Stage IV (4 days--2 weeks): recovery if liver regeneration is adequate; complete histologic recovery is possible.
\section*{Alternative Names}
APAP Level; Paracetamol Level; Tylenol Level; Acetaminophen Serum Concentration; Acetaminophen Overdose Screen
\section*{Principle}
Acetaminophen is measured by enzymatic assay or immunoassay, with results reported in mcg/mL (mg/L). The revised Rumack--Matthew nomogram is used only for a single acute ingestion with a reasonably known time, using a concentration drawn at least 4 hours after ingestion. In the United States and Canada, the treatment line begins at 150 mcg/mL (mg/L) at 4 hours; other countries may use different protocols. The historical 200-line must not be presented as the current treatment threshold. The nomogram is not valid for repeated supratherapeutic ingestion or an unknown time of ingestion, and extended-release products or delayed absorption may require serial concentrations and toxicology advice.
\section*{History}
Acetaminophen was first synthesized in 1878 by Harmon Northrop Morse and was introduced into clinical medicine in the 1950s as a safer alternative to phenacetin and acetanilide, which caused methemoglobinemia. It was marketed in the United States as Tylenol in 1955 and in the United Kingdom as Panadol in 1956 \cite{prescott2000paracetamol}. The hepatotoxicity of acetaminophen overdose was first recognized in 1966, when Davidson and Eastham reported fulminant hepatic failure after massive ingestion \cite{davidson1966acute}. Barry Rumack and Henry Matthew then developed the nomogram relating the serum level to time after ingestion, published in \textit{Pediatrics} in 1975 \cite{rumack1975acetaminophen}. N-acetylcysteine was introduced as the antidote following studies by Prescott and colleagues in the late 1970s \cite{prescott1979intravenous} and remains the standard of care.
\section*{Why This Test Is Ordered}
Evaluation and management of acute acetaminophen overdose (intentional or accidental). A serum acetaminophen level must be drawn 4 hours after ingestion (or as soon as possible thereafter). Chronic acetaminophen overdose in patients at risk (malnutrition, chronic alcoholism --- both deplete glutathione and induce CYP2E1, increasing NAPQI production at lower acetaminophen doses).
\section*{Primary vs Secondary Test}
This is a secondary, therapeutic-monitoring test rather than a diagnostic test. It is drawn to assess the risk of hepatotoxicity and to guide N-acetylcysteine therapy after a known or suspected acetaminophen overdose.
\section*{How This Test Is Performed}
Blood is collected by standard venipuncture into a plain serum tube (red-top) or a serum separator tube (SST, gold-top); plasma may also be used. The sample is drawn at least 4 hours after ingestion (or as soon as possible thereafter) so that absorption is complete. The serum acetaminophen concentration is measured by an enzymatic assay (acetaminophen amidohydrolase) or by particle-enhanced turbidimetric inhibition immunoassay (PETINIA) on automated analyzers. Results are reported in mcg/mL and plotted on the Rumack-Matthew nomogram against the time since ingestion; results are typically available within 1--2 hours for urgent overdose management.
\section*{What Preparation Is Needed}
No fasting or special preparation is required. The critical variable is timing: the level must be drawn at least 4 hours after ingestion of an immediate-release formulation, and if the ingestion time is uncertain or absorption may be delayed, a repeat level is drawn 4 hours after the first. The dose should be held while the level is pending in the setting of suspected overdose.
\section*{Contraindications and Precautions}
\begin{itemize}
  \item Interpret the level against the Rumack-Matthew nomogram only for a single acute ingestion of immediate-release acetaminophen; the nomogram is not valid for chronic excessive dosing, extended-release products, or ingestions of unknown timing
  \item Do not withhold N-acetylcysteine while awaiting the level if the history suggests a significant overdose; it can be started empirically
  \item Levels drawn before 4 hours after ingestion are not interpretable on the nomogram and must be repeated
  \item Co-ingestion of opioids, anticholinergic drugs, or ethanol delays gastric emptying and absorption, shifting the toxicokinetic curve to the right; repeat the level 4 hours after the first
  \item Chronic alcoholism, malnutrition, and fasting deplete glutathione and induce CYP2E1, increasing susceptibility to hepatotoxicity at lower acetaminophen doses
  \item Check for co-ingested salicylates in the same sample; pregnancy does not contraindicate N-acetylcysteine
\end{itemize}
\section*{Risks}
Minimal risk. Slight pain or bruising at the venipuncture site. Rarely: hematoma, infection, or vasovagal syncope.
\section*{What Happens After the Test}
The acetaminophen level is plotted against the applicable treatment line only when the nomogram is valid. N-acetylcysteine is started promptly when the concentration is at or above the applicable treatment line, when the ingestion time is unknown or delayed and toxicity cannot be excluded, or when clinical and laboratory findings support toxicity. Treatment duration is individualized; acetaminophen concentration, AST/ALT, INR, renal function, and clinical status are followed serially with poison-center or toxicology input.
\section*{Understanding Results}
An acetaminophen concentration at or above the applicable treatment line requires prompt NAC under the local toxicology protocol. NAC is most effective when started early but remains beneficial after hepatotoxicity develops. Co-ingestion of opioids or anticholinergic drugs can delay absorption; extended-release products and delayed or staggered ingestions may require repeat concentrations and cannot be interpreted with a single nomogram point.
\section*{Normal Results}
There is no universal therapeutic serum range for routine dosing that should be used to manage overdose. Toxicity is interpreted using the applicable nomogram only for a valid single acute ingestion. NAC should be considered urgently for an unknown-time or repeated ingestion with a detectable concentration, abnormal AST/ALT, or other evidence of toxicity, in consultation with a poison center or toxicologist.
\section*{Follow-up Tests}
AST, ALT, total bilirubin, PT/INR, serum creatinine/BUN, serum glucose (hypoglycemia in Stage III hepatic failure), arterial blood gas (lactic acidosis), serum phosphate (hyperphosphatemia, a marker of poor prognosis in acetaminophen-induced fulminant hepatic failure) \cite{larson2005acetaminophen}, and King's College Criteria for liver transplantation.
\section*{References}
\bibliographystyle{unsrtnat}
\bibliography{references}
\clearpage
\section*{Regulatory Status and Authorization Date}
This chapter describes a test category, not a specific commercial product. FDA, CDSCO, EU IVDR, and MHRA approval or registration dates therefore cannot be assigned without the assay manufacturer, product name, instrument, and intended use. For laboratory-developed tests, an individual agency approval date may not exist. See the Preface for the interpretation of regulatory status.

\clearpage
